\documentclass[11pt,a4paper]{article}
\usepackage[T1]{fontenc}
\usepackage[utf8]{inputenc}
\usepackage{lmodern,microtype}
\usepackage[margin=.85in]{geometry}
\usepackage{amsmath,amssymb,amsthm,mathtools,aliascnt}
\usepackage{booktabs,tabularx,array,enumitem,needspace,float}
\usepackage[section]{placeins}
\usepackage{xcolor,tikz,mdframed}
\usetikzlibrary{arrows.meta,positioning,calc,decorations.pathreplacing}
\usepackage[hidelinks]{hyperref}
\usepackage[nameinlink,noabbrev]{cleveref}
\usepackage{fancyhdr}
\hypersetup{pdftitle={Sync HotStuff RR v4.3},pdfsubject={Regional finality and L1-authorized rotation}}
\setlength{\headheight}{14pt}
\pagestyle{fancy}\fancyhf{}
\lhead{\small Sync HotStuff RR v4.3}\rhead{\small L1-authorized regional sessions}
\cfoot{\thepage}
\setlength{\parindent}{0pt}\setlength{\parskip}{4pt}
\setlength{\emergencystretch}{2em}
\renewcommand{\topfraction}{.92}
\renewcommand{\textfraction}{.06}
\renewcommand{\floatpagefraction}{.85}
\makeatletter
\setlength{\@fptop}{0pt}\setlength{\@fpbot}{0pt plus 1fil}
\makeatother
\setlist{nosep,leftmargin=*}
\newcolumntype{Y}{>{\raggedright\arraybackslash}X}
\definecolor{rrRed}{HTML}{A51C30}
\definecolor{rrBlue}{HTML}{2563EB}
\definecolor{rrTeal}{HTML}{0F766E}
\newcommand{\RR}[1]{{\color{rrRed}#1}}
\newcommand{\QC}{\mathcal C}
\newcommand{\FC}{\mathcal F}
\newcommand{\Hash}{\mathsf H}
\newcommand{\rank}{\operatorname{rank}}
\newcommand{\End}{\operatorname{End}}
\newcommand{\TB}{\mathit{TB}}
\newcommand{\tagg}[1]{\ensuremath{\mathsf{#1}}}
\newtheorem{definition}{Definition}[section]
\newaliascnt{assumption}{definition}
\newtheorem{assumption}[assumption]{Assumption}
\aliascntresetthe{assumption}
\crefname{assumption}{assumption}{assumptions}
\Crefname{assumption}{Assumption}{Assumptions}
\newcounter{protocol}
% Protocol boxes float to the top of a page, so long boxes leave no gaps.
\newenvironment{protocol}[2]{\begin{figure}[htbp]\refstepcounter{protocol}\label{#1}
\begin{mdframed}[linewidth=.6pt,skipabove=0pt,skipbelow=0pt,
innertopmargin=7pt,innerbottommargin=7pt]\small
\textbf{Protocol \theprotocol. #2}\par\smallskip}{\end{mdframed}\end{figure}}
\crefname{protocol}{Protocol}{Protocols}
\Crefname{protocol}{Protocol}{Protocols}
\title{Sync HotStuff RR v4.3\\\large Regional Finality and L1-Authorized Rotation}
\author{}\date{12 October 2026}
\begin{document}
\maketitle
\begin{abstract}
This draft extends Sync HotStuff RR v3.3~\cite{rr33} to several regions
that take turns. One turn of one region is a \emph{session}. L1 fixes
each session's stake, leader schedule, and quota of DATA blocks. Inside
a session, the region runs v3.3: proposal timers commit blocks, and
terminal certificates rotate leaders. A session ends in one of two ways.
In a regular handoff, the region commits its last quota block, and the
next live region continues from that block. In a forced handoff, an
honest multisignature committee posts a force instruction on L1. The
instruction names the last DATA height that the external chain keeps
from the session. The old region may still commit higher blocks, but
the external chain drops them. The force is an external hard fork,
effective when its activation L1 block is finalized; that block supplies
the activation timestamp. Both handoffs use the same bootstrap
steps. A recovery that opens at the quota limit proposes one DRAIN
block. It carries no data; it only commits the last DATA block. This
draft specifies the model and the protocols. Security and complexity
analyses are deferred.
\end{abstract}

\section{Overview}
\label{sec:overview}
RR v3.3~\cite{rr33} revises standard-synchronous Sync HotStuff~\cite{sync}
for one region. This draft lets several regions take turns producing
blocks. Session $r$ is one turn, served by region $i_r$. The design has
three layers.
\begin{itemize}
\item \textbf{L1} is the public source of authority. It fixes each
session's voters, leaders, and quota. It also carries status updates
and force instructions (\cref{sec:model}).
\item \textbf{The regional engine} runs v3.3 inside one session, with a
quota on DATA blocks (\cref{sec:regional,sec:protocols}). Its replicas
are the members of region $i_r$.
\item \textbf{The external chain} is the chain that clients and later
regions accept. Every node builds it from the DATA prefixes that the
sessions retain. The next region starts from its end
(\cref{sec:external,sec:xprotocols}). Nodes are the replicas of all
regions and any other verifier.
\end{itemize}
A backbone carries blocks and receipts between regions. It never
chooses the next region; L1 records and regional certificates do.

A session ends by a \emph{regular} handoff, after its last quota block
$T_r$ commits, or by a \emph{forced} handoff, when an L1 force
instruction cuts its external history at a DATA height $h_*$.
\Cref{fig:endings} in \cref{sec:examples} shows both endings.

A view is $(r,g,t)$: session $r$, recovery generation $g$, and leader
term $t$. Views are ordered lexicographically. Each session starts at
$(r,0,0)$. \Cref{tab:views} lists the four ways to change view. The
first two happen inside region $i_r$, as in v3.3. The members of the
next region take the last two.

\begin{table}[htbp]\centering\small
\caption{The four view changes.}
\label{tab:views}
\begin{tabularx}{\linewidth}{@{}>{\raggedright\arraybackslash}p{.19\linewidth}>{\raggedright\arraybackslash}p{.12\linewidth}Y>{\raggedright\arraybackslash}p{.23\linewidth}@{}}
\toprule
Change & New view & Trigger & Commit timers\\
\midrule
Normal rotation & $(r,g,t+1)$ & A terminal QC of term $t$, for DATA below $H_r$ & Keep running.\\
Recovery & $(r,g+1,0)$ & A quit proof for generation $g$, then a $\Delta$ wait & Canceled in generation $g$.\\
Regular handoff & $(r+1,0,0)$ & Ready on a REGULAR witness, which holds a receipt for $T_r$ (\cref{sec:ready}) & The old engine is already complete.\\
Forced handoff & $(r+1,0,0)$ & Ready on a FORCE witness, which holds an activated force & The old engine is unchanged.\\
\bottomrule
\end{tabularx}
\end{table}
A QC for DATA at $H_r$, or for a DRAIN, starts no new term. A status
update without a force changes no view. Recovery keeps the session's
quota.

\section{Model and assumptions}
\label{sec:model}

\subsection{L1 control state}
\label{sec:l1}
At each finalized L1 height, nodes read the same authority records and
status-update list (\cref{tab:l1}). The authority
$\tagg{AUTH}_{r,g}$ gives the configuration, quota, and leader schedule
for generation $(r,g)$ on the branch opened by the session boundary.
Finalized L1 block headers also give their authenticated timestamps.
The multisignature committee posts status updates. A regular handoff
needs no new L1 transaction.

\begin{table}[htbp]\centering\small
\caption{L1 records used by the protocol.}
\label{tab:l1}
\begin{tabularx}{\linewidth}{@{}p{.27\linewidth}Y@{}}
\toprule
Field & Meaning\\
\midrule
\multicolumn{2}{@{}l}{\textbf{Authority} $\tagg{AUTH}_{r,g}$}\\
Stake configuration & The members $V_i$ of each region, their public
keys and voting weights, and protocol parameters such as $K$ and
$\Delta_i$.\\
Leader schedule & The leaders within each region. Once a boundary
chooses $i_r$, this gives $L(r,g,t)\in V_{i_r}$.\\
Regional quota $\TB_r$ & The number of new DATA blocks that session $r$
may add. Recovery does not reset it.\\
\midrule
\multicolumn{2}{@{}l}{\textbf{Status update}}\\
Status map $\chi$ & Live or off, for each region.\\
Force flag & True if the update forces a handoff.\\
Target session $r$ & The session to cut, identified by its index and
opening boundary. Used only if the flag is true.\\
Anchor height $h_*$ & The last DATA height that the external chain keeps
from session $r$. Used only if the flag is true. It is a height, not a
block hash.\\
Activation height $\lambda_a$ & The update takes effect once L1 block
$\lambda_a$ is finalized. Its activation timestamp $\theta_a$ is that
block's L1 timestamp, not a local message-arrival time.\\
\bottomrule
\end{tabularx}
\end{table}

The handoff chooses region $i_r$ (\cref{sec:witness}). The finalized
authority supplies the voters and leader schedule within that region.

Status updates follow these rules.
\begin{itemize}
\item Each update is posted at least a configured notice period before
its activation height. Nodes may prepare during this period. The update
has no effect before activation.
\item Activation heights increase from one update to the next. Nodes
apply updates in L1 order. The \emph{effective status map}
$\chi(\lambda)$ at L1 height $\lambda$ is the map of the latest update
whose activation height is at most $\lambda$.
\item A force update also marks its target region off. An update whose
force flag is false changes only the status map.
\item Nodes check each update's authorization and L1 finality. A
finalized record never changes.
\end{itemize}
The initial configuration fixes the cyclic region order $\mathcal L_R$,
the initial status map, and the genesis block. The first live region in
$\mathcal L_R$ serves session $0$. If no region is live, initialization
waits.

\subsection{Stake and faults}
\label{sec:stake}
Member $v$ of region $i$ has voting weight $\omega_{i,v}>0$. Let
$W_i=\sum_{v\in V_i}\omega_{i,v}$. Byzantine members of region $i$ hold
total weight at most $F_i<W_i/2$. Every regional certificate needs
matching signatures from distinct members with total weight greater than
$F_i$. So every certificate has a correct signer, and the correct members
alone can form one. This replaces the $f+1$ of $n=2f+1$ in v3.3.

Membership and weights stay fixed for a whole session, including its
recoveries. Slashing and status updates do not change them. A later
session checks old certificates against the old snapshot. The bound
$F_i$ counts every key of the snapshot that is compromised while its
certificates are still accepted. Compromise of old keys after that is
outside this model.

Correct replicas keep their protocol state and committed blocks.
Signatures cannot be forged, and hashes are binding.

\subsection{Assumptions}
\label{sec:assumptions}
\begin{assumption}[Finalized L1 records]
\label{ass:authority}
At every finalized L1 block height, the authority records and the ordered
list of status updates through that height are common knowledge. These
finalized records never change.
\end{assumption}
A replica uses only finalized records. If an authority needed for entry
is not yet finalized, the replica waits for it.

\begin{assumption}[Authenticated synchronous communication and global clock]
\label{ass:network}
\leavevmode\par
\begin{itemize}
\item Messages use authenticated communication channels.
\item Within region $i$, the actual delay of a correct-to-correct
message is $\delta\le\Delta_i$; $\delta$ may differ across messages.
\item All nodes share a global clock.
\end{itemize}
\end{assumption}
The delay bound includes delivery and processing of complete validation
dependencies. Each replica records its own entry and voting events on
the global clock.

\begin{assumption}[Cross-region backbone]
\label{ass:backbone}
A backbone eventually delivers each session's committed blocks, with
their regional receipts, to all other nodes. Its delay has no fixed
bound.
\end{assumption}

\begin{assumption}[Honest multisignature committee]
\label{ass:committee}
The multisignature committee is honest. Every update it posts follows
the rules of \cref{sec:l1,sec:force}.
\end{assumption}

\subsection{Conditions for progress}
\label{sec:progress}
The deferred liveness analysis will also use these conditions.
\begin{itemize}
\item \textbf{Schedules.} Correct leaders get repeated turns in views
that are actually entered. This includes recovery initializers and
DATA-producing turns across sessions. Turns in views that are never
reached do not count. 
\item \textbf{Exclusion.} Leaders exposed for equivocation are eventually
excluded by later authorities, as in v3.3's accountability assumption.
\item \textbf{Committee service.} The committee gives working regions
enough uninterrupted time to produce DATA and to complete handoffs.
\item \textbf{Requests.} A correct client's request that remains
admissible eventually reaches every region that continues to serve.
Each request fits in one DATA batch. Each node keeps pending requests
in the order it first receives them, with only finitely many requests
ahead of any one request (\cref{sec:requests}).
\end{itemize}
A correct leader takes the oldest admissible requests in its own queue
that fit and are not already in its proposed ancestry. New arrivals do
not move ahead of requests already queued. Different nodes need not
have the same queue order. If no request can be included, the leader
may propose an empty DATA block.

\section{Regional engine: definitions}
\label{sec:regional}

\Cref{tab:notation} collects the notation used below. Anchor indices
are always written in full.
\begin{table}[htbp]\centering\small
\caption{Notation.}\label{tab:notation}
\begin{tabularx}{\linewidth}{@{}p{.29\linewidth}Y@{}}
\toprule Symbol & Meaning\\\midrule
$(r,g,t)$; $j$ & Session, recovery generation, leader term; position
within that term. Indices refer to one opening branch.\\
$i_r$; $V_i$ & Region serving session $r$; members of region $i$.\\
$\omega_{i,v}$; $W_i$; $F_i$ & Member weight; total regional weight;
upper bound on Byzantine weight, with $F_i<W_i/2$.\\
$\Delta_i$; $K$ & Regional message-delay bound; maximum DATA blocks
in a leader term.\\
$\tagg{AUTH}_{r,g}$; $\mathit{cfg}_r$ & Finalized authority for the
generation; identifier of the session's fixed voting snapshot.\\
$A_{r,g,t}$ & Anchor immediately before term $(r,g,t)$.\\
$A_{r,0,0}$; $A_{r,g,0}$ & Inherited session anchor; generation anchor
chosen at initialization or recovery.\\
$h(B)$; $k_r(B)$ & DATA height from genesis; physical position within
session $r$, counting DATA and DRAIN.\\
$\TB_r$; $H_r$; $T_r$ & Session DATA quota; last allowed DATA height;
DATA block at that height.\\
$m_{r,g,t}$ & Number of fresh positions in the term.\\
$L(r,g,t)$ & Leader of the view.\\
$\QC_{r,g,t}(B)$; ANCHOR QC & Proposal certificate for $B$;
certificate re-certifying $A_{r,g,0}$.\\
$\FC_r(B)$ & Regional receipt: matching FINAL signatures of weight
$>F_{i_r}$, validated by \cref{def:receipt}.\\
$Q_r^0$; $Q_v$; $Q^*$ & Inherited anchor's receipt; replica's entry
lock; justification selected by the recovery leader.\\
$e_v(r,g,t)$; $\tau_v(P)$; $a_v$ & Local view-entry time; time of first
validated vote for proposal $P$; time of receiving a quit proof.\\
$\rho$; $\lambda(B)$ & Block's finalized L1 reference; its L1 height.\\
$\chi(\lambda)$; $\mathcal L_R$ & Effective live/off map at L1 height
$\lambda$; cyclic region order.\\
$h_*$; $\lambda_a$; $\theta_a$ & Forced retained DATA height;
activation L1 height; timestamp of that L1 block.\\
$\lambda_s$; $\lambda_r^{\rm open}$ & Successor-selection L1 height;
selection height of the boundary that opened session $r$.\\
$\End(r)$; $\mathcal W$; $\Xi$ & Retained endpoint; boundary witness;
its finalized L1 proofs.\\
$E$; $\tagg{Ready}_v(\mathcal W)$ & Entry proof; completion of the
successor's boundary, history, and authority validation.\\
\bottomrule
\end{tabularx}
\end{table}

\subsection{Blocks, positions, and L1 references}
\label{sec:blocks}
The anchor $A_{r,g,t}$ is the block immediately before the first fresh
position of view $(r,g,t)$. We always write all three indices.
\begin{itemize}
\item $A_{r,0,0}$ is the DATA block inherited from the preceding
session; $A_{0,0,0}$ is genesis. It keeps the regional receipt from the
session that produced it. Genesis has a synthetic receipt.
\item For $g>0$, NEWVIEW chooses the recovery anchor $A_{r,g,0}$
(\cref{alg:recovery}). It is certified but may be uncommitted.
\item For $t>0$, $A_{r,g,t}$ is the terminal DATA block of term
$(r,g,t-1)$, certified by that term's terminal QC.
\end{itemize}
The anchor of a view never changes. The generation starts at
$A_{r,g,0}$; later terms use their own anchors $A_{r,g,t}$. ANCHOR voting
re-certifies only $A_{r,g,0}$, once per generation.

Each block has two positions.
\begin{itemize}
\item The \emph{DATA height} $h(B)$ counts the DATA blocks on the chain
from genesis to $B$.
\item The \emph{physical position} $k_r(B)$ counts DATA and DRAIN blocks
within session $r$, with $k_r(A_{r,0,0})=0$. This local count does not change
the signed fields of $A_{r,0,0}$.
\end{itemize}
A new block of session $r$ has the form
\[
 B=(r,k,h,\mathit{type},x,\Hash(B_{\rm parent}),\rho),
 \qquad \mathit{type}\in\{\tagg{DATA},\tagg{DRAIN}\},
\]
with $k=k_r(B_{\rm parent})+1$. A DATA block carries a batch $x$ of
requests and has $h=h(B_{\rm parent})+1$. A DRAIN block carries no batch
and has $h=h(B_{\rm parent})$.

\paragraph{L1 references.}
The reference $\rho$ names a finalized L1 block. It holds that L1
block's height $\lambda(B)$ and hash, and a registry proof for the L1
control records up to that height. These rules hold:
\begin{itemize}
\item along DATA ancestry, $\lambda$ never decreases;
\item the first DATA block of session $r$ has $\lambda\ge
\lambda^{\rm open}_r$, the selection height of the boundary that opened
session $r$ (\cref{sec:witness});
\item a DRAIN block copies the reference of $T_r$.
\end{itemize}
The reference of $T_r$ gives the starting L1 height for selecting the
regular successor (\cref{sec:witness}). It cannot protect a block from
an activated force, even if the reference predates the force.

Regional validity uses only the session's fixed authority. Later L1
updates do not change it. A regionally valid block can still be left out
of the external chain (\cref{sec:external}).

\subsection{Session quota, term positions, and DRAIN}
\label{sec:quota}
\paragraph{1. Fix the session endpoint.}
Session $r$ may add $\TB_r\ge1$ DATA blocks after $A_{r,0,0}$. Its last
DATA height is
\[
 H_r=h(A_{r,0,0})+\TB_r.
\]
The DATA block at this height is $T_r$. Recovery never changes $H_r$.
The quota counts the DATA blocks after $A_{r,0,0}$ in a proposal's
ancestry, whether committed or not. DRAIN, resends, and ANCHOR votes
use no quota.

\paragraph{2. Produce DATA while space remains.}
If $h(A_{r,g,0})<H_r$, generation $(r,g)$ is a \emph{DATA generation}.
Each term begins at its anchor $A_{r,g,t}$ and has
\[
 m_{r,g,t}=\min\{K,\;H_r-h(A_{r,g,t})\}
\]
new positions, numbered $j=0,\ldots,m_{r,g,t}-1$. At position $j$, the
DATA block has
\[
 h=h(A_{r,g,t})+j+1,\qquad
 k=k_r(A_{r,g,t})+j+1.
\]
A new position is called a \emph{fresh position}; its block is a
\emph{fresh block}. The anchor itself is not fresh. The last position
$j=m_{r,g,t}-1$ is the \emph{terminal position}, and its QC is the
\emph{terminal QC}.

If the terminal DATA height is below $H_r$, its QC admits term
$(r,g,t+1)$, with that terminal block as $A_{r,g,t+1}$. If it is $H_r$,
there is no next term. Only the last DATA term can have fewer than $K$
positions. The initial generation is always a DATA generation, since
$\TB_r\ge1$.

\paragraph{3. Use DRAIN when recovery starts at the endpoint.}
If $g>0$ and $h(A_{r,g,0})=H_r$, generation $(r,g)$ is a
\emph{drain generation}. It has only term $t=0$ and position $j=0$,
with $m_{r,g,0}=1$. Its one fresh block is DRAIN, whose parent is
$A_{r,g,0}$ and whose justification is the generation's ANCHOR QC.
It carries no requests, keeps DATA height $H_r$, and has no next term.

DRAIN supplies a commit timer. ANCHOR voting starts no timer, so
re-certifying $A_{r,g,0}$ alone does not commit it. At height $H_r$, a
DATA child would exceed the quota. The DRAIN child's timer instead
commits its ancestors, including $T_r$.
\begin{itemize}
\item A DATA generation that reaches $H_r$ does not append DRAIN.
It waits for $T_r$'s timer or recovers after a quit.
\item A later drain generation may start from an earlier DRAIN.
Each interrupted recovery may add one DRAIN, so their total number
is not bounded.
\item The choice depends on the recovery anchor's DATA height, not on
whether a replica has already committed $T_r$.
\end{itemize}
DRAIN stays in the regional log. The external chain contains only DATA,
and the next session inherits $T_r$.

\paragraph{Example.}
Let $h(A_{r,0,0})=100$, $\TB_r=10$, and $K=4$, so $H_r=110$.
\begin{center}\small
\begin{tabular}{@{}lll@{}}\toprule
View & Anchor height & Fresh DATA heights\\\midrule
$(r,0,0)$ & 100 & 101--104\\
$(r,0,1)$ & 104 & 105--108\\
$(r,0,2)$ & 108 & 109--110\\\bottomrule
\end{tabular}
\end{center}
Recovery from DATA 107 starts a new generation whose first term produces
108--110. Recovery from certified DATA 110 starts a drain generation.
In both cases, the endpoint remains 110.

\subsection{Certificates and signed messages}
\label{sec:certificates}
\Cref{tab:certificates} lists the certificates of a session. Each needs
matching signatures from distinct members of $V_{i_r}$ with total weight
$>F_{i_r}$.

\begin{table}[htbp]\centering\small
\caption{Certificates of a session.}
\label{tab:certificates}
\begin{tabularx}{\linewidth}{@{}>{\raggedright\arraybackslash}p{.25\linewidth}>{\raggedright\arraybackslash}p{.24\linewidth}Y@{}}
\toprule
Certificate & Signed message & Purpose\\
\midrule
Proposal QC $\QC_{r,g,t}(B)$ & DATA-VOTE or DRAIN-VOTE & Justifies a
child. A terminal QC below $H_r$ also admits the next term.\\
ANCHOR QC & ANCHOR-VOTE & Re-certifies the opening $A_{r,g,0}$ in
generation $g$, before any fresh proposal.\\
BLAME proof & BLAME & Freezes a generation (\cref{sec:quit}).\\
Regional receipt $\FC_r(B)$ & FINAL & Proves that a correct replica
committed DATA block $B$.\\
\bottomrule
\end{tabularx}
\end{table}

A correct replica casts at most one fresh vote per $(r,g,t,j)$ and one
ANCHOR vote per $(r,g)$. DATA and DRAIN votes follow the same rule but
are different signed types. No QC and no ANCHOR vote starts a commit
timer.

\paragraph{Signed context.}
Every signature binds the session's opening boundary, the voting
configuration, session $r$, region $i_r$, and the message type. This
separates ordinary continuation from a later force fork
(\cref{sec:witness}). Messages of a view also bind the L1
authority and $(r,g,t)$. The other signed fields are:
\begin{itemize}
\item \textbf{PROPOSE and fresh votes:} the opening, position $j$, block
type, physical position, DATA height, and block hash;
\item \textbf{NEWVIEW and ANCHOR-VOTE:} the identity of the chosen
opening. NEWVIEW also binds its justification and entry proof;
\item \textbf{STATUS:} the sender's entry lock and entry proof. No
opening is needed;
\item \textbf{BLAME:} the view only. It binds no opening, block hash, or
progress count. All complaints for a view thus sign the same statement,
even before NEWVIEW;
\item \textbf{FINAL:} session-stable fields only (below).
\end{itemize}
Message displays omit these context fields.

\paragraph{Rank and the base justification.}
Within session $r$, a QC formed in term $(r,g,t)$ for block $B$ has rank
$(g,t,k_r(B))$, compared lexicographically. The ANCHOR QC of generation
$g$ has rank $(g,0,k_r(A_{r,g,0}))$, just below the generation's first
fresh QC. A fixed total order on certified content breaks ties. Signer
sets and encodings do not matter. An opening's identity is
$(g,k_r(A_{r,g,0}),h(A_{r,g,0}),\Hash(A_{r,g,0}))$.

The anchor $A_{r,0,0}$ also has a \emph{base justification} $Q_r^0$: the
receipt of $A_{r,0,0}$ in the boundary witness that opened session $r$
(\cref{sec:witness}). $Q_r^0$ ranks below every QC of session $r$.
Recovery uses it when no QC of the session is known. A QC of another
session is never a justification in session $r$, so it cannot replace
$A_{r,0,0}$.

\paragraph{Regional receipts.}
After a replica commits a DATA block $B$ of session $r$, by $B$'s own
timer or as an ancestor, it sends
\[
 \langle\tagg{FINAL},\mathit{cfg}_r,r,h(B),\Hash(B)\rangle_v
\]
to its region and to the backbone. Matching FINAL messages of weight
$>F_{i_r}$ form the regional receipt $\FC_r(B)$. The receipt carries $B$ and the ancestry needed to validate it, or
refers to a validated local cache.

\begin{definition}[Valid regional receipt]
\label{def:receipt}
A regional receipt $\FC_r(B)$ is valid when all of the following hold:
\begin{enumerate}[label=(\roman*)]
\item $B$ is a DATA block of session $r$, extends $A_{r,0,0}$, and has
$h(A_{r,0,0})<h(B)\le H_r$;
\item distinct members in session $r$'s fixed voting snapshot have
valid signatures on the same
$\langle\tagg{FINAL},\mathit{cfg}_r,r,h(B),\Hash(B)\rangle$,
with total weight $>F_{i_r}$ and the signed context specified above; and
\item $B$ and its required ancestry pass the block-validation rules of
\cref{sec:admission}.
\end{enumerate}
The verifier checks this evidence; it does not replay the signers'
timers. The weight threshold guarantees at least one correct signer,
who sends FINAL only after committing $B$.
\end{definition}
Here $\mathit{cfg}_r$ identifies session $r$'s fixed voting snapshot.
When we write $\FC(B)$ without a session subscript, we mean the receipt
from the session that produced $B$. The inherited anchor keeps that
receipt; it does not acquire a new one.
\begin{itemize}
\item FINAL binds no generation or term. So FINAL messages for one block
combine across recoveries. The block hash keeps the block's signed
context.
\item If a DRAIN commits $T_r$, replicas send FINAL for $T_r$, never for
the DRAIN.
\item Replicas of session $r$ send no FINAL for $A_{r,0,0}$. It keeps its
receipt from session $r-1$.
\item A replica may commit the DATA prefix of a valid receipt
(\cref{def:receipt}) of its own session. Receipts of other sessions serve only as history or anchor
evidence. They never change this session's locks or log.
\end{itemize}

\subsection{Admission, validation, and engine states}
\label{sec:admission}
\paragraph{Admission.}
Replica $v$ admits views as follows.
\begin{itemize}
\item $(r,0,0)$ is admitted at $v$'s Ready event for the boundary that
opens session $r$ (\cref{def:ready}). The witness is the entry proof.
The NEWVIEW must choose exactly $A_{r,0,0}$.
\item $(r,g,0)$ with $g>0$ is admitted once $v$ holds
$\tagg{AUTH}_{r,g}$ and a quit proof for generation $g-1$. That proof is
the entry proof. The NEWVIEW is justified by a QC of an earlier
generation of session $r$, or by $Q_r^0$.
\item $(r,g,t)$ with $t>0$ is admitted once $v$ also holds the terminal
QCs of terms $0,\ldots,t-1$, each for DATA below $H_r$.
\end{itemize}
Term zero is admitted before its NEWVIEW arrives. Its progress checks
therefore run even if the initial leader is silent. Fresh proposals need
the generation's ANCHOR QC. This holds in every generation, including
generation $0$ of session $0$; unlike v3.3, genesis has no shortcut. A
message is \emph{admissible} if its view is admitted and its signatures
are valid.

\paragraph{Complete validation.}
Validation checks signatures, authority, admission, block type, quota,
both positions, the L1-reference rules, the parent rule
(\cref{tab:parent}), parent QCs, hash links, and application validity
throughout the required ancestry. Every relevant ancestor includes its
signed proposal and block data; a chain of certificate hashes alone is
not enough. Dependencies come from the message or a validated local
cache. Correct senders send complete validation bundles. A reference is
allowed only if the receiver can resolve it within the delay bound of
\cref{ass:network}. Bundles list each dependency once.

\paragraph{Evidence takes priority.}
Before validating a message or discarding a duplicate, inspect all signed
material in it for a quit proof (\cref{sec:quit}). A proof needs its own
valid signatures and admission witnesses. A valid proof counts even if
the message that carries it is invalid or incomplete. Compare incoming
material with all retained material, not only with objects inside the
message.

\paragraph{One proposal-processing rule.}
Every admitted fresh proposal of the current generation that $v$ has not
voted for invokes the Vote step. This includes proposals in attached
histories, proposals of earlier terms, and proposals whose blocks are
already committed as ancestors. A proposal's identity is
$(r,g,t,j,\Hash(B))$. The signer set or encoding of its parent QC does
not change it. Deduplicate by this identity after evidence inspection.
Old-generation proposals are retained without new votes or timers.
Future-generation messages wait for entry.

\paragraph{Engine states.}
A replica runs one engine for each session opening of its region.
Engines on different branches remain separate, even if they use the
same session index. The engine of
session $r$ is in one of three states.
\begin{itemize}
\item \textbf{Active:} the current generation is admitted and unfrozen.
The engine proposes, votes, and changes terms.
\item \textbf{Frozen:} the current generation has a quit proof. The
engine stops proposals, votes, and term changes. It still processes
evidence, old votes, QCs, receipts, and recovery.
\item \textbf{Complete:} the replica holds $\FC_r(T_r)$ satisfying
\cref{def:receipt}.
Production in session $r$ is over, even if the external handoff is still
pending. The engine only serves data.
\end{itemize}
Committing $T_r$ locally does not complete the engine. Other members may
still need this replica's recovery votes before a receipt forms. Each
message goes to the engine identified by its session opening, so
evidence from one session or branch never freezes another. These states are separate from the external
states of \cref{sec:states}. The external chain may move to a later
session while an old engine still runs.

\begin{table}[htbp]\centering\small
\caption{Parent selection for the next fresh position $j$ of term $(r,g,t)$.}
\label{tab:parent}
\begin{tabularx}{\linewidth}{@{}l Y Y@{}}
\toprule
Condition & $B_{\mathrm{parent}}$ & $Q_{\mathrm{parent}}$\\
\midrule
$j=0,t=0$ (DATA or DRAIN) & Opening $A_{r,g,0}$ & Its ANCHOR QC in $(r,g,0)$\\
$j=0,t>0$ & Anchor $A_{r,g,t}$, the terminal block of $(r,g,t-1)$ & Its proposal QC in $(r,g,t-1)$\\
$1\le j<m_{r,g,t}$ & Position $j-1$ of $(r,g,t)$ & Its proposal QC in $(r,g,t)$\\
\bottomrule
\end{tabularx}
\end{table}

\section{Regional protocols}
\label{sec:protocols}
Protocols 1 and 2 run inside the engine of one session $r$. They are the
protocols of v3.3 with the changes listed below. \RR{Red text} marks
every difference from v3.3, in the boxes, in the message-processing
order, and in the quit proofs. Explanations stay outside the boxes.

\paragraph{Notation.}
All anchors and view times use their full session indices. The displayed
regional messages omit $r$, but their signatures bind the session's
opening boundary and $r$. In these two protocols, $\Delta$ and $F$ mean
$\Delta_{i_r}$ and $F_{i_r}$.

\paragraph{Changes from v3.3.}
\begin{enumerate}[label=\arabic*.]
\item \emph{Stake.} Certificates need weight $>F$, not $f+1$ signatures.
\item \emph{Quota.} A term has $m_{r,g,t}\le K$ positions. No term follows
$H_r$.
\item \emph{DRAIN.} A drain generation proposes one DRAIN. Its vote
counts as progress.
\item \emph{Receipts.} Replicas send FINAL after committing DATA. A
receipt for $T_r$ completes the engine (new step~5 of \cref{alg:steady}).
\item \emph{Timers.} Completion also cancels timers. A force never does.
\item \emph{Recovery entry.} After the $\Delta$ wait, a replica enters
when $\tagg{AUTH}_{r,g+1}$ is finalized. $Q_r^0$ can serve as the entry
lock and as the justification.
\item \emph{Sessions.} Each message goes to the engine of its session.
Session $0$ opens by bootstrap (\cref{alg:bootstrap}).
\end{enumerate}

\paragraph{Message-processing order.}\leavevmode
\RR{Replica $v$ first passes each message to the engine of the session
that the message binds. L1 updates, witnesses, and receipts of other
sessions go to \cref{alg:external}.} The engine handles each received
message in this order:
\begin{enumerate}[label=(\arabic*)]
\item check all signed material for a quit proof;
\item check completeness and validity;
\item process attached DATA \RR{and DRAIN} proposals under Vote, parents
first;
\item process attached terminal QCs under step~4 of \cref{alg:steady},
in increasing term order;
\item process the message itself.
\end{enumerate}
When a message and a timer fall at the same time, $v$ handles the
message first.

\subsection{Steady state}
\label{sec:steady}
\Cref{alg:steady} gives the steady state of a session: proposals, votes,
timers, term rotation, and completion.

\begin{protocol}{alg:steady}{Propose, vote, commit, normal rotation\RR{, and completion}}
Replica $v$ runs these steps \RR{in the engine of session $r$}. Its
current term is $(r,g,t)$.
\begin{enumerate}[label=\textbf{\arabic*)},itemsep=3pt]
\item \textbf{Propose.} Run this step only if $v=L(r,g,t)$, $g$ is
unfrozen\RR{, and session $r$ is not complete}.
\RR{In a DATA generation,} for each position
$j=0,\ldots,\RR{m_{r,g,t}}-1$, in order:
 \begin{itemize}
 \item Wait until $v$ holds $B_{\mathrm{parent}}$ and
       $Q_{\mathrm{parent}}$ for $j$ (\cref{tab:parent}).
 \item Build \RR{$B=(r,k,h,\tagg{DATA},x,\Hash(B_{\mathrm{parent}}),\rho)$,
       with $k$ and $h$ from \cref{sec:quota} and an L1 reference $\rho$}.
       The batch $x$ follows \cref{sec:progress}.
 \item Broadcast
       $P=\langle\tagg{PROPOSE},g,t,j,B,Q_{\mathrm{parent}}\rangle_v$
       with its validation bundle.
 \item Process $P$ under Vote.
 \end{itemize}
 \RR{In a drain generation, propose once, at $(t,j)=(0,0)$. The block is
 a DRAIN with parent $A_{r,g,0}$, no batch, DATA height $H_r$, and the L1
 reference of $T_r$. Broadcast and process it as above.}
 Do not wait for the parent to commit. Never sign two different
 blocks for one position. A resend repeats the same $P$.

\item \textbf{Vote.} Run this step on every DATA \RR{or DRAIN} proposal
$P$ of term $(r,g_P,t_P)$ at position $j_P$. Here $t_P$ may be earlier
than $t$. The step also applies to proposals inside attached histories.
 \begin{enumerate}[label=(\alph*)]
 \item If $P$ or its attached material contains a quit proof for $g$,
       run Quit and stop.
 \item If $g_P<g$, retain $P$ and stop. If $g_P>g$, buffer $P$
       until $v$ enters generation $g_P$.
 \item If $g$ is frozen \RR{or session $r$ is complete}, retain $P$ and
       stop.
 \item If a dependency is missing, retain $P$, request the dependency,
       and stop. Rerun this step when it arrives.
 \item If $P$ is invalid, discard it and stop. \RR{Validity includes the
       block type, the quota, both positions, and the L1 reference.}
 \item If $v$ has already voted for $P$, stop.
       A duplicate never restarts a timer.
 \item Otherwise, let $\tau_v(P)$ be the current time. Forward $P$ and
       its bundle to all replicas. Broadcast
       $\langle\tagg{DATA\mbox{-}VOTE},g,t_P,j_P,\Hash(B)\rangle_v$
       \RR{(DRAIN-VOTE if $B$ is a DRAIN)}.
       Start the commit timer of $P$. It expires at $\tau_v(P)+2\Delta$.
 \end{enumerate}

\item \textbf{Commit.} When the timer of $P$ expires, commit $B$ and
its uncommitted ancestors. \RR{For each DATA block $D$ of session $r$
that this commits, send
$\langle\tagg{FINAL},\mathit{cfg}_r,r,h(D),\Hash(D)\rangle_v$ to the region
and the backbone. A committed DRAIN yields FINAL for $T_r$, never for
itself. Relay each receipt of session $r$ that $v$ obtains to the region
and the backbone.} Only Quit \RR{and Completion} cancel a timer. Entering a
new term does not cancel it. A QC never starts a timer. \RR{A force
instruction never cancels a timer.}

\item \textbf{Terminal QC and rotation.} Run this step the first time
$v$ obtains a valid terminal QC for a term $(r,g,t')$.
``Obtains'' means that $v$ assembles the QC from votes or finds it in
any message, including as another proposal's $Q_{\mathrm{parent}}$.
 \begin{enumerate}[label=(\alph*)]
 \item Process each unseen proposal in the QC's history under Vote,
       parents first.
 \item If $g$ is now frozen, stop. \RR{If the QC certifies DATA at
       $H_r$, or a DRAIN, stop: no next term exists.}
 \item Send the QC and its validation bundle to $L(r,g,t'+1)$.
 \item If $t\le t'$, enter $(r,g,t'+1)$ and set $e_v(r,g,t'+1)$ to the
       current time. Cancel the progress checks of all earlier
       terms.
 \end{enumerate}
 A QC never moves $v$ to an earlier term. It never changes a recorded
 entry time.

\item \RR{\textbf{Completion.} Run this step the first time $v$ obtains a
regional receipt $\FC_r(T_r)$ satisfying \cref{def:receipt}.}
 \begin{itemize}\color{rrRed}
 \item Commit $T_r$ and its uncommitted ancestors.
 \item Mark session $r$ complete. Stop its proposals, votes, progress
       checks, term changes, and recovery. Cancel its unexpired timers.
 \item Relay $\FC_r(T_r)$ and its bundle to the region and the backbone.
 \item Keep serving data. \Cref{alg:external} runs the handoff.
 \end{itemize}
 \RR{A force instruction never runs this step.}
\end{enumerate}
\end{protocol}

The acceptance time $\tau_v(P)$ is when complete validation and the first
vote occur, not an earlier arrival of an incomplete packet. Each replica
records its own entry time for each term. In session $r$, $e_v(r,0,0)$ is
set at Ready, in step~2 of \cref{alg:bootstrap}. For $g>0$, $e_v(r,g,0)$ is
set in step~3 of \cref{alg:recovery}. For $t>0$, $e_v(r,g,t)$ is set in
step~4 of \cref{alg:steady}. A term's progress checks therefore start
when the term is admitted locally, not when its leader's first proposal
arrives. A silent new leader cannot delay them.

A proposal or QC at height $H_r$ does not end the session. Without a
quit, the timer of $T_r$ commits it; no later block is needed. Only a
receipt for $T_r$ completes the engine.

\subsection{Quit proofs and recovery}
\label{sec:quit}
A quit proof freezes a generation. \Cref{alg:recovery} then moves the
session to the next generation.

\paragraph{Complete quit proofs.}
For an admitted generation $g$ \RR{of session $r$}, a proof consists of
either:
\begin{enumerate}[label=(\alph*)]
\item distinct signatures \RR{of total weight $>F$} on the same
$\langle\tagg{BLAME},g,t\rangle$ for an admitted term, with that term's
admission witness;
\item two different admissible DATA \RR{or DRAIN} blocks signed by
$L(r,g,t)$ at the same $(r,g,t,j)$, with their admission \RR{and quota}
witnesses; or
\item two valid NEWVIEWs signed by $L(r,g,0)$ selecting different opening
blocks for generation $g$.
\end{enumerate}
\RR{A DATA block beyond the quota, or a DRAIN outside a drain generation,
is invalid. It is rejected and never counts for (b).}
Different openings count even when their blocks are comparable, because
they fix different starting heights $h(A_{r,g,0})$. The height formula then
maps the same position to different heights.

Different QCs for the same opening block are not equivocation. Evidence
from any admitted term can freeze its generation, including a term not
yet locally reached. Old-generation proofs are archival and may be
submitted to L1. \RR{Evidence from another session never freezes this
engine.} Future-generation proofs are buffered until entry. Neither
resets an unrelated recovery.

\begin{protocol}{alg:recovery}{Blame, quit, and recovery}
Replica $v$ runs these steps \RR{in the engine of session $r$}. Its
current term is $(r,g,t)$.
\begin{enumerate}[label=\textbf{\arabic*)},itemsep=3pt]
\item \textbf{Blame.} \RR{Let $q=m_{r,g,t}$ in a DATA generation and
$q=1$ in a drain generation.} For each $p=1,\ldots,\RR{q}$, run a
progress check at time
\[
 e_v(r,g,t)+(2p+4)\Delta.
\]
\RR{Schedule the check for $p=1$ at entry, even if the opening is not yet
known. Schedule the others once $q$ is known. A check whose time has
already passed runs at once.}
At that time, broadcast $\langle\tagg{BLAME},g,t\rangle_v$ if all of
these hold:
 \begin{itemize}
 \item $v$'s current term is still $(r,g,t)$, $g$ is unfrozen\RR{, and
       session $r$ is not complete};
 \item $v$ has cast DATA \RR{or DRAIN} votes for fewer than $p$ positions
       of term $(r,g,t)$ (ANCHOR votes do not count);
 \item $v$ has not yet sent BLAME for $(r,g,t)$.
 \end{itemize}

\item \textbf{Quit.} Run this step the first time $v$ holds a quit proof
for its current generation $g$.
 \begin{itemize}
 \item Freeze $g$. Stop all progress checks, proposals, votes, and
       term changes in $g$. Cancel every unexpired commit timer of $g$.
 \item Set $a_v$ to the current time. Broadcast the proof.
 \item Keep all commitments. Keep collecting votes and QCs of $g$.
       They can raise the entry lock chosen in step~3. \RR{Keep
       processing regional receipts (step~5 of \cref{alg:steady}).}
 \item If the proof is of type (b) or (c), submit it to L1. Do
       not wait for L1 to process it.
 \end{itemize}

\item \textbf{Status and entry.} \RR{Skip this step if session $r$ is
complete.} \RR{Wait until at least $a_v+\Delta$ and until
$\tagg{AUTH}_{r,g+1}$ is finalized. Set $e_v(r,g+1,0)$ to the time both
conditions hold, and read that authority from L1.} Then:
 \begin{itemize}
 \item Set the entry lock $Q_v$ to the highest-ranked QC
       (DATA\RR{, DRAIN,} or ANCHOR) that $v$ knows from generations
       $0,\ldots,g$\RR{, or to $Q_r^0$ if it knows none}.
       $Q_v$ never changes after this. QCs that arrive later do not
       update it.
 \item Enter $(r,g+1,0)$ as unfrozen and start its progress checks.
       \RR{Keep $A_{r,0,0}$, $H_r$, and the voting snapshot of session $r$.}
 \item Send $\langle\tagg{STATUS},g+1,0,Q_v,E\rangle_v$ and its
       validation bundle to $L(r,g+1,0)$. Here $E$ is the quit proof for
       $g$, which is also the entry proof of $g+1$.
 \item Process any buffered proofs for $g+1$.
 \end{itemize}

\item \textbf{New-view.} Run this step only at $L(r,g+1,0)$. Wait $2\Delta$
after its own entry. If $g+1$ is still current and unfrozen:
 \begin{itemize}
 \item Let $Q^*$ be the highest-ranked QC from generations $0,\ldots,g$
       that the leader knows, including the QCs in STATUS
       messages\RR{; if it knows none, let $Q^*=Q_r^0$}.
 \item Let $A_{r,g+1,0}$ be the block that $Q^*$ certifies. \RR{If
       $h(A_{r,g+1,0})<H_r$, generation $g+1$ is a DATA generation.
       Otherwise it is a drain generation.}
 \item Broadcast
       $\langle\tagg{NEWVIEW},g+1,0,A_{r,g+1,0},Q^*,E\rangle_{L(r,g+1,0)}$
       with its validation bundle. Choose only one opening block for
       $g+1$.
 \end{itemize}

\item \textbf{Anchor vote.} On receiving a NEWVIEW for $g+1$:
 \begin{enumerate}[label=(\alph*)]
 \item If it contains a quit proof for $g+1$, run Quit and stop.
 \item Validate its leader, entry proof $E$, opening block $A_{r,g+1,0}$,
       $Q^*$, and dependencies. If data are missing, request them and
       wait without voting. If the message is invalid, discard it.
 \item Vote only if all of these hold: $g+1$ is current and unfrozen; $v$
       has cast no ANCHOR vote in $g+1$; and $\rank(Q^*)\ge\rank(Q_v)$.
 \item To vote, forward the NEWVIEW to all replicas and broadcast
       $\langle\tagg{ANCHOR\mbox{-}VOTE},g+1,0,\Hash(A_{r,g+1,0})\rangle_v$.
       Start no timer. Use no fresh position.
 \end{enumerate}
 Matching ANCHOR votes \RR{of total weight $>F$} form the ANCHOR QC.
 When $L(r,g+1,0)$ obtains it, the leader starts Propose in
 \cref{alg:steady} at $j=0$\RR{; in a drain generation, this is the
 DRAIN}. Every replica, including one that cast no ANCHOR vote,
 processes DATA \RR{and DRAIN} proposals backed by a valid ANCHOR QC
 under Vote.
\end{enumerate}
\end{protocol}

The entry-lock test governs ANCHOR voting only. A replica that has
committed $T_r$ locally still takes part in recovery until it holds
$\FC_r(T_r)$, because other members may need its votes.

\section{External chain and handoff}
\label{sec:external}

\subsection{Session states and the external chain}
\label{sec:states}
Every node keeps each session in one of three external states
(\cref{fig:states}).
\begin{itemize}
\item \textbf{Open:} its endpoint is not fixed. DATA blocks with valid
regional receipts, up to $H_r$, are provisional output.
\item \textbf{Truncated:} an activated force has fixed the retained
endpoint $\End(r)=h_*$, but the anchor data or successor may still be
missing.
\item \textbf{Sealed:} a boundary witness fixes the retained anchor
$A_{r+1,0,0}$, successor $i_{r+1}$, and selection height $\lambda_s$.
The next session can start when its members are Ready.
\end{itemize}
A regular boundary keeps $\End(r)=H_r$ and
$A_{r+1,0,0}=T_r$. A force boundary keeps $\End(r)=h_*$ and uses the DATA
block at $h_*$, or $A_{r,0,0}$ itself if $h_*=h(A_{r,0,0})$.

Sealing fixes a boundary on its branch. An activated force can replace
a regular boundary and discard every later session descended from it.
Ordinary receipts and status changes cannot do this. Once a force fixes
the endpoint of its target session, regular output never restores the
removed suffix. The current external session is the first unsealed
session on the branch selected by the finalized L1 force rules.

\begin{figure}[htbp]\centering
% External states of session r (Section 5.1).
% Open -> Sealed (regular handoff); Open -> Truncated -> Sealed (forced handoff).
% A force can replace an already sealed regular boundary.
\begin{tikzpicture}[x=1cm,y=1cm,font=\fontsize{9}{10.5}\selectfont,
  st/.style={draw,rounded corners=3pt,minimum width=2.9cm,minimum height=1cm,
    align=center,inner sep=3pt},
  tr/.style={-{Stealth[length=1.8mm]},semithick},
  lbl/.style={font=\fontsize{8.5}{10}\selectfont,align=center,inner sep=2pt}]
  \node[st,draw=black!60,fill=black!4] (open) at (0,0)
    {\textbf{Open}\\[-1pt]{\fontsize{8}{9}\selectfont output is provisional}};
  \node[st,draw=rrRed,fill=rrRed!6] (trunc) at (4.7,-1.9)
    {\textbf{Truncated}\\[-1pt]{\fontsize{8}{9}\selectfont $\End(r)=h_*$ is fixed}};
  \node[st,draw=rrTeal,fill=rrTeal!10] (sealed) at (11.2,0)
    {\textbf{Sealed}\\[-1pt]{\fontsize{8}{9}\selectfont anchor and successor fixed}};
  \draw[tr] (open) -- node[lbl,above] {regular receipt and live successor; $\End(r)=H_r$} (sealed);
  \draw[tr] (open.south) |- node[lbl,pos=.74,above] {force activates} (trunc.west);
  \draw[tr] (trunc.east) -| node[lbl,pos=.26,above] {anchor data and a live
    successor} (sealed.south);
  \draw[tr,draw=rrRed] (sealed.north) to[bend right=28]
    node[lbl,above,text=rrRed] {force replaces a regular boundary}
    (trunc.north);
\end{tikzpicture}

\caption{External states of session $r$. A regular boundary seals the
session. An activated force truncates it, including when a normal
boundary was already sealed, and removes the discarded descendants.
Regional engine states are separate (\cref{sec:admission}).}
\label{fig:states}
\end{figure}

\paragraph{External acceptance.}
A DATA block of a truncated or sealed session belongs to the external
chain only if it has a valid regional receipt (\cref{def:receipt}),
extends $A_{r,0,0}$, and lies in the retained DATA prefix through
$\End(r)$. The session must itself be on the branch selected by the
activated L1 forces. Genesis and these prefixes, in session order,
form the external chain. DRAIN blocks never belong to it.

A regional receipt proves regional commitment. External finality also
applies the L1 force rule: a block on a discarded branch is rejected,
even when its regional receipt remains valid. In this model, normal
sealing is therefore subject to a later authorized force hard fork.

\subsection{Force instructions and their activation time}
\label{sec:force}
A force update names a target session $r$ and its opening boundary, an
anchor height $h_*$, and an activation L1 height $\lambda_a$. Define
\[
 \theta_a=\operatorname{Time}_{\mathrm{L1}}(\lambda_a),
\]
the timestamp in the activation L1 block. The update has no effect
before that block is finalized. Once it is finalized, the force is
active. Neither the posting time nor a node's message-arrival time is
its activation time.

The honest committee posts a force only when:
\begin{itemize}
\item the target is the session it is replacing, identified by its
opening boundary on the currently retained branch, and no earlier
force has fixed that target's endpoint;
\item $h(A_{r,0,0})\le h_*\le H_r$, preserving the inherited prefix; and
\item the chosen anchor has a valid regional receipt and its data and
ancestry are available outside the region. The committee chooses the
highest such DATA block from the prefix produced while the region
worked correctly. If $h_*=h(A_{r,0,0})$, the inherited receipt suffices.
\end{itemize}
The update specifies a height, not a hash. The receipt identifies the
block at that height, using regional safety on the target branch.

\paragraph{The external hard-fork rule.}
When the activation block is finalized, every node keeps the target
session only through $h_*$. Blocks above $h_*$ are discarded whether
they were proposed or committed before or after $\theta_a$. The force
replaces any regular handoff of the target. All descendants of that old
boundary are discarded too, even if $h_*=H_r$. This rule also applies
to a node that learns the force later.
The node reconstructs the same forced branch from finalized L1 records;
its earlier local handoff does not defeat the force.

A force leaves the old regional BFT unchanged: it does not cancel timers,
change locks or quota, erase commitments, or supply a quit proof.
Old-branch blocks and receipts remain regional evidence, but cannot
restore discarded external output. No proof that the old engine stopped
is required. If several force records name the same target, the first
one to activate fixes its endpoint; later ones change only the status
map. A force whose target was removed by an earlier force also changes
only the status map. The activation block and its timestamp identify
the force.

\subsection{Boundary witness and successor}
\label{sec:witness}
\paragraph{Successor.}
Let $\tagg{NextLive}(i,\chi)$ be the first region after $i$ in
$\mathcal L_R$, other than $i$, that is live in $\chi$. The search covers
at most one cycle. If no such region exists,
$\tagg{NextLive}(i,\chi)=\bot$. Let
\[
 \lambda_0=\begin{cases}
 \lambda(T_r) & \text{for a regular handoff,}\\
 \max\{\lambda_a,\lambda(A_{r+1,0,0})\} & \text{for a forced handoff.}
 \end{cases}
\]
The \emph{selection height} $\lambda_s$ is the least L1 height
$\lambda\ge\lambda_0$ with $\tagg{NextLive}(i_r,\chi(\lambda))\ne\bot$.
The successor is $i_{r+1}=\tagg{NextLive}(i_r,\chi(\lambda_s))$, and
$\lambda^{\rm open}_{r+1}=\lambda_s$. So every node computes the same
successor. If no other region is live at $\lambda_0$, the boundary waits
for the first later status map that makes one live. Once chosen, the
successor never changes. For session $0$, $\lambda^{\rm open}_0$ is the
initial L1 height.

\paragraph{Regular and force boundaries.}
A regular witness can be formed from $\FC_r(T_r)$ and the selected live
successor. It posts nothing to L1. A force witness instead uses the
activated force and the receipt at $h_*$. Once a force is active for the
target, its witness takes precedence over every regular witness for that
target, including one formed earlier. Message-arrival order is irrelevant.

\paragraph{The witness.}
A boundary witness has the form
\[
 \mathcal W=(\mathit{mode},r,i_{r+1},A_{r+1,0,0},
 \FC(A_{r+1,0,0}),\lambda_s,\Xi),
 \quad \mathit{mode}\in\{\tagg{REGULAR},\tagg{FORCE}\}.
\]
The proof $\Xi$ supplies the finalized L1 authority and status-update
records needed to check the boundary, including the map at $\lambda_s$.
For FORCE, it also proves the activation block, its timestamp
$\theta_a$, and the target's opening boundary.
\begin{itemize}
\item \textbf{REGULAR:} $A_{r+1,0,0}=T_r$, with a valid receipt and
exactly $\TB_r$ DATA blocks after $A_{r,0,0}$.
\item \textbf{FORCE:} $A_{r+1,0,0}$ is the DATA block at $h_*$, with its
original receipt. If $h_*=h(A_{r,0,0})$, the failed region signs nothing
new.
\end{itemize}
A regular witness on a branch removed by an active force remains historical
evidence; it no longer authorizes external continuation. Any node can
validate a current witness and seal its boundary. Only the successor's
members use it to start an engine.

\paragraph{Branch identity.}
A session is identified by its opening boundary as well as its index
$r$. The boundary identity uses the preceding session's opening boundary,
its mode, anchor hash, successor, and selection height. For FORCE it
also identifies the activated L1 update. Different certificate encodings or signer sets do not change it.
Every regional signed message binds this identity. Each branch uses
the finalized authority for its session and generation; the branch
identity does not require an extra L1 record. If a force replaces
a regular handoff, the new session $r+1$ is on a different branch, even
if an old session used the same numeric index. Its engine, votes, locks,
and timers are separate. Old-branch evidence cannot freeze it or justify
its proposals.

Session $0$ uses a synthetic witness: $A_{0,0,0}$ is genesis with a
synthetic receipt, and $i_0$ is the first live region in $\mathcal L_R$.

\subsection{The Ready event}
\label{sec:ready}
\begin{definition}[Ready]
\label{def:ready}
Let $v$ be a member of region $i'$. $\tagg{Ready}_v(\mathcal W)$ occurs
the first time $v$ has completely validated all of the following:
\begin{enumerate}[label=(\roman*)]
\item a sealed witness $\mathcal W$ on the current external branch
that names $i'$ as the successor of session $r$;
\item the anchor and any application history that $v$ is missing; and
\item the finalized $\tagg{AUTH}_{r+1,0}$.
\end{enumerate}
A REGULAR witness must show $\FC_r(T_r)$ and the full quota. A FORCE
witness must instead show the finalized activation block and its
timestamp, together with the retained cutoff. Ready does not occur if
$v$ has already entered this session or a later session on the same
branch. Entry on a discarded branch does not prevent entry on the
forced branch. An incomplete witness stays pending
until its data and authority arrive.
\end{definition}
Ready is the only event that lets a region start the next session. A
block at height $H_r$, its QC, an announced force, or a newer status map
is not enough. Ready starts the local entry clock. Remote proposal
times, private QC assembly, and L1 installation times do not. Members of
$i'$ process boundary bundles before entry. Other messages of session
$r+1$ wait until entry.

\subsection{Pending requests}
\label{sec:requests}
Every node keeps a queue of pending client requests. A request is
\emph{admissible} when it passes the application's admission checks.
The node places it at the back when first received; receiving it again
does not change its place. Correct leaders use their own queue order,
as specified in \cref{sec:progress}.

A request leaves the queue when included in a sealed retained prefix.
Regional commitment alone does not remove it. If a force discards a
prefix that had been sealed by a regular handoff, requests removed by that prefix
become pending again at their original queue positions unless they
occur in the retained history. Nodes keep enough request history to
restore them, and leaders include them
like any other pending request.

\section{Cross-region protocols}
\label{sec:xprotocols}
Every node runs \cref{alg:external}. Members of the successor region
also run \cref{alg:bootstrap}. Both protocols are new relative to v3.3.
\Cref{alg:external} is all red. \Cref{alg:bootstrap} reuses steps~3--5
of v3.3's recovery protocol; red marks its differences from those steps.

\begin{protocol}{alg:external}{Watch L1, apply forces, and seal boundaries}
\color{rrRed}
Every node $u$ runs these steps for the current external branch.
\begin{enumerate}[label=\textbf{\arabic*)},itemsep=3pt]
\item \textbf{Watch.} Continuously read finalized L1 authority and
status-update records. Receive regional blocks and receipts from the
backbone, and fetch missing data. Keep the proofs and ancestry needed
for a boundary. While a session is open, its receipted DATA output is
provisional.

\item \textbf{Apply a force.} When L1 block $\lambda_a$ is finalized,
apply the force identified by that block and timestamp $\theta_a$.
\begin{itemize}
\item Locate the target session using its opening boundary. If an
earlier force has removed that target or already fixed its endpoint,
apply only the status map.
\item Otherwise set $\End(r)=h_*$ and mark the target truncated, even
if its regular boundary was already sealed.
\item Discard its output above $h_*$ and every session descended from
the replaced regular boundary. Restore their requests to the pending
queue when absent from the retained history. Keep their regional logs
as old-branch evidence.
\item Make the target the current external session on the forced
branch. Fetch the anchor at $h_*$, its receipt, and its ancestry.
\item Compute the selection height and successor. When the data,
successor, and matching authority are available, form a FORCE witness
and run Seal.
\end{itemize}
The cutoff holds while data or a successor is missing. A newly received
old regular witness does not undo it.

\item \textbf{Regular completion.} When a valid $\FC_r(T_r)$ arrives for
the current session, first apply all forces active in the finalized L1
records known at that height. If the session has a forced endpoint,
retain the receipt only as regional evidence. Otherwise:
\begin{itemize}
\item Check the receipt and that $T_r$'s ancestry contains exactly
$\TB_r$ DATA blocks after $A_{r,0,0}$.
\item Compute the selection height and successor. If no other region
is live, keep $T_r$ as a candidate and wait.
\item Form a REGULAR witness and run Seal. A force that activates later
can still replace this regular boundary under step 2.
\end{itemize}

\item \textbf{Seal.} Validate the witness and its branch against the
finalized L1 records. A regular witness for a discarded branch is not
accepted. Then:
\begin{itemize}
\item Record the retained endpoint, anchor $A_{r+1,0,0}$, successor
$i_{r+1}$, and selection height $\lambda_s$.
\item Append the retained DATA prefix and remove its requests from the
pending queue.
\item Mark session $r$ sealed on this branch and move to $r+1$.
\item If $u\in V_{i_{r+1}}$, run \cref{alg:bootstrap} for this opening
boundary.
\end{itemize}
\end{enumerate}
\end{protocol}

New L1 records, blocks, receipts, or authority rerun the applicable steps,
even while a regional engine is frozen. Repeating a step for the same
boundary has no further effect. Nodes with the same finalized L1 records
and regional evidence choose the same branch and boundary. Before entry,
the successor applies all status updates effective at $\lambda_s$.
Session $0$ starts from Seal with its synthetic witness.



\begin{protocol}{alg:bootstrap}{Bootstrap the successor region}
Replica $v$ \RR{of region $i_{r+1}$} runs these steps \RR{for the sealed
witness $\mathcal W$ of session $r$}. Here $\Delta$ and $F$ refer to
region $i_{r+1}$.
\begin{enumerate}[label=\textbf{\arabic*)},itemsep=3pt]
\item \RR{\textbf{Ready.} On $\tagg{Ready}_v(\mathcal W)$
(\cref{def:ready}), relay $\mathcal W$ and its complete bundle to all
members of $i_{r+1}$. A receiver's validation of this bundle is its own
Ready event. It runs the next steps at once. It does not wait for the
backbone or for a leader message.}

\item \textbf{Status and entry.} \RR{Install session $r+1$ with anchor
$A_{r+1,0,0}$, authority $\tagg{AUTH}_{r+1,0}$, quota $\TB_{r+1}$, and
opening L1 height $\lambda_s$.} Then:
 \begin{itemize}
 \item Set the entry lock $Q_v$ to \RR{the base justification
       $Q^0_{r+1}$}.
 \item Enter \RR{$(r+1,0,0)$} as unfrozen and start its progress checks.
       \RR{Set $e_v(r+1,0,0)$ once.}
 \item Send $\langle\tagg{STATUS},\RR{r+1,0,0},Q_v,E\rangle_v$ and its
       validation bundle to $L(r+1,0,0)$. Here \RR{$E=\mathcal W$} is
       the entry proof.
 \item Process any buffered proofs for \RR{generation $0$ of session
       $r+1$}.
 \end{itemize}
 \RR{A duplicate of the same witness never resets the entry time, the
 quota, or the lock. A force witness that replaces a regular boundary
 starts a separate branch engine.}

\item \textbf{New-view.} Run this step only at $L(r+1,0,0)$. Wait
$2\Delta$ after its own entry. If \RR{$(r+1,0,0)$} is still current and
unfrozen, broadcast
$\langle\tagg{NEWVIEW},\RR{r+1,0,0,A_{r+1,0,0},Q^0_{r+1},\mathcal W}\rangle_{L(r+1,0,0)}$
with its validation bundle. \RR{The opening must be exactly $A_{r+1,0,0}$.}

\item \textbf{Anchor vote.} Run step~5 of \cref{alg:recovery} for
\RR{generation $0$ of session $r+1$, with $E=\mathcal W$}. When the
leader obtains the ANCHOR QC, it starts Propose in \cref{alg:steady} at
$j=0$\RR{, with DATA height $h(A_{r+1,0,0})+1$ and physical position $1$.
Later failures in session $r+1$ use \cref{alg:recovery}.}
\end{enumerate}
\end{protocol}

\paragraph{Bootstrap timing.}
Let $e$ be the first Ready time at a correct member of $i_{r+1}$. The
relay in step~1 makes every correct member enter by $e+\Delta$. Assume a
correct initial leader and no quit. The leader sends NEWVIEW by
$e+3\Delta$. The ANCHOR QC reaches it by $e+5\Delta$. The first DATA
proposal reaches every correct member by $e+6\Delta$. This is no later
than each member's first progress check, at
$e_v(r+1,0,0)+6\Delta\ge e+6\Delta$. The backbone delay before Ready has
no fixed bound.

\section{Examples}
\label{sec:examples}
\Cref{fig:endings} uses the running example: $h(A_{r,0,0})=100$, $\TB_r=10$,
and $K=4$, so $H_r=110$.

\paragraph{Regular ending after a DRAIN.}
Generation $0$ has three terms: DATA 101--104, 105--108, and 109--110.
The QC of DATA 110 starts no next term. If the timers of DATA 110 expire,
it commits, and its receipt completes the session. Suppose instead that
a quit cancels these timers first. Generation $1$ then opens on certified
DATA 110, so it is a drain generation. Its only fresh block is one DRAIN.
The DRAIN's timer commits DATA 110 as an ancestor. Replicas send FINAL
for DATA 110, not for the DRAIN. Another quit would need one more
generation and one more DRAIN. In every case, session $r+1$ starts with
DATA 111, whose parent is DATA 110.

\paragraph{Forced ending.}
The committee posts a force for session $r$ with $h_*=106$ and
$\lambda_a=1200$. Its activation timestamp is the timestamp of L1
block 1200. The old region may still commit DATA 107 and 108,
before or after activation. Once L1 block 1200 is finalized, every node
keeps session $r$ only through DATA 106. It drops output above 106, including any descendants of a normal
handoff from that output, and restores their requests when absent
from the retained prefix. After validating DATA 106
and its receipt, the successor opens $(r+1,0,0)$. It proposes its own
DATA 107, whose parent is DATA 106. The old regional log does not
change. Its receipts for 107 and 108 never bring those blocks back.

\begin{figure}[htbp]\centering
% Two ways a session ends (Section 7), with h(A_r)=100, TB_r=10, K=4.
% (a) Regular ending after one DRAIN.  (b) Forced ending at h_*=106.
% Teal: external chain. Gray: regional log only. Blue: session r+1.
% Arrows point from parent to child.
\begin{tikzpicture}[x=1cm,y=1cm,font=\fontsize{9}{10.5}\selectfont,
  blk/.style={draw,rounded corners=1.5pt,minimum height=.62cm,
    minimum width=.95cm,inner sep=2pt},
  ext/.style={blk,draw=rrTeal,fill=rrTeal!12},
  reg/.style={blk,draw=black!45,fill=black!6,text=black!65},
  nxt/.style={blk,draw=rrBlue,fill=rrBlue!10},
  link/.style={-{Stealth[length=1.4mm]},semithick,draw=black!70},
  lane/.style={anchor=east,font=\fontsize{8}{9}\selectfont,text=black!70},
  tag/.style={font=\fontsize{8}{9}\selectfont,inner sep=1pt},
  ttl/.style={anchor=west,font=\bfseries\fontsize{9}{10.5}\selectfont},
  note/.style={anchor=north west,align=flush left,text width=5.7cm,
    inner sep=0pt,font=\fontsize{8.5}{10}\selectfont},
  brace/.style={decorate,decoration={brace,amplitude=3pt},draw=black!55}]
  % ---------------- (a) regular ending ----------------
  \node[ttl] at (-1.55,1.3) {(a) Regular ending after one DRAIN};
  \node[lane] at (.35,0) {session $r$};
  \node[lane] at (.35,-1.2) {session $r{+}1$};
  \node[ext] (a100) at (.95,0) {100};
  \node[ext,minimum width=1.5cm] (a101) at (2.55,0) {101--104};
  \node[ext,minimum width=1.5cm] (a105) at (4.35,0) {105--108};
  \node[ext] (a109) at (5.8,0) {109};
  \node[ext] (a110) at (7.0,0) {110};
  \node[reg,dashed] (aD) at (8.35,0) {DRAIN};
  \node[nxt] (a111) at (8.35,-1.2) {111};
  \draw[link] (a100)--(a101); \draw[link] (a101)--(a105);
  \draw[link] (a105)--(a109); \draw[link] (a109)--(a110);
  \draw[link] (a110)--(aD);
  \draw[link] (a110.south)--(a111.west);
  \node[tag,below=2pt of a100] {$A_{r,0,0}$};
  \node[tag,anchor=north east] at ([xshift=-1pt,yshift=-2pt]a110.south west) {$T_r$};
  \node[tag,above=3pt of a101] {term 0};
  \node[tag,above=3pt of a105] {term 1};
  \draw[brace] ([yshift=3pt]a109.north west)--([yshift=3pt]a110.north east)
    node[tag,midway,above=4pt] {term 2};
  \node[tag,above=3pt of aD] {gen.\ 1};
  \node[note] at (9.4,.5) {A quit cancels the timers of 110 before it
    commits. Generation~1 opens on 110, so it is a drain generation.
    Its one DRAIN's timer commits 110, and FINAL is sent for 110 only.
    Session $r{+}1$ starts at 111, with parent 110.};
  % ---------------- (b) forced ending ----------------
  \def\yb{-3.65}\def\yn{-4.85}
  \node[ttl] at (-1.55,\yb+1.3) {(b) Forced ending at $h_*=106$};
  \node[lane] at (.35,\yb) {session $r$};
  \node[lane] at (.35,\yn) {session $r{+}1$};
  \node[ext] (b100) at (.95,\yb) {100};
  \node[ext,minimum width=1.5cm] (b101) at (2.55,\yb) {101--104};
  \node[ext] (b105) at (4.0,\yb) {105};
  \node[ext] (b106) at (5.2,\yb) {106};
  \node[reg] (b107) at (6.65,\yb) {107};
  \node[reg] (b108) at (7.85,\yb) {108};
  \node[nxt] (n107) at (6.65,\yn) {107};
  \draw[link] (b100)--(b101); \draw[link] (b101)--(b105);
  \draw[link] (b105)--(b106); \draw[link] (b106)--(b107);
  \draw[link] (b107)--(b108);
  \draw[link] (b106.south)--(n107.west);
  \draw[rrRed,thick,densely dashed] (5.92,\yb+.55)--(5.92,\yb-.5);
  \node[tag,text=rrRed,above] at (5.92,\yb+.55) {cutoff};
  \node[tag,below=2pt of b100] {$A_{r,0,0}$};
  \node[tag,anchor=north east] at ([xshift=-1pt,yshift=-2pt]b106.south west) {$h_*$};
  \node[note] at (9.4,\yb+.5) {The force activates when L1 block 1200 is finalized, using its timestamp.
    The external chain keeps session $r$ only through 106. The old region
    may still commit 107 and 108; the external chain drops them and
    any normal-handoff descendants.
    Session $r{+}1$ proposes its own DATA 107, with parent 106.};
  % ---------------- legend ----------------
  \def\yl{-5.95}
  \node[ext,minimum width=.5cm,minimum height=.34cm] at (-.9,\yl) {};
  \node[anchor=west,tag] at (-.55,\yl) {external chain};
  \node[reg,minimum width=.5cm,minimum height=.34cm] at (2.3,\yl) {};
  \node[anchor=west,tag] at (2.65,\yl) {regional log only};
  \node[nxt,minimum width=.5cm,minimum height=.34cm] at (5.7,\yl) {};
  \node[anchor=west,tag] at (6.05,\yl) {session $r{+}1$};
\end{tikzpicture}

\caption{The two ways a session ends, with $h(A_{r,0,0})=100$, $\TB_r=10$,
and $K=4$. Arrows point from parent to child. (a) A quit leaves DATA 110
certified but uncommitted; one DRAIN commits it. (b) A force cuts the
external chain at 106; the old region's 107 and 108 stay in its log
only.}
\label{fig:endings}
\end{figure}

\begingroup\small
\begin{thebibliography}{9}
\bibitem{rr33}
Sync HotStuff RR v3.3: Proposal Timers and Terminal-QC Handoff.
Specification draft, 12 October 2026 (\texttt{sync\_hotstuff\_rr\_v3\_3.tex}).
\bibitem{sync}
I. Abraham, D. Malkhi, K. Nayak, L. Ren, and M. Yin.
\emph{Sync HotStuff: Simple and Practical Synchronous State Machine Replication}.
IEEE Symposium on Security and Privacy, 2020, pp. 106--118.
\href{https://conferences.computer.org/sp/pdfs/sp/2020/349700a106.pdf}{Published paper}.
\end{thebibliography}
\endgroup
\end{document}
